\documentclass{svproc}
\usepackage{url}
\def\UrlFont{\rmfamily}
\usepackage{amsmath}
\usepackage{array}
\usepackage{graphicx}

\begin{document}
\mainmatter 

\title{NLP Assignment: Part 1 Answers}
\titlerunning{Part 1 Answers} 
\author{Top Indian Student}
\authorrunning{Top Indian Student}
\institute{University}

\maketitle 

\begin{abstract}
This document contains the step-by-step solutions for Part 1 of the NLP assignment, which involves working out a toy Information Retrieval (IR) system focusing on vector space models, TF-IDF representations, cosine similarity, and word sense ambiguity.
\keywords{Information Retrieval, TF-IDF, Vector Space Model, Boolean Retrieval}
\end{abstract}

\section*{Part 1: Working out a Toy IR System}

\subsection*{1. Preprocessing and Inverted Index}
Here are the original documents:
$d_1$: The star in our solar system provides heat and light. \\
$d_2$: That Hollywood star walked the red carpet for the movie premiere. \\
$d_3$: Astronomers observe distant stars and galaxies using telescopes.

\textbf{Step 1: Tokenization} \\
Converting sentences to lowercase and stripping punctuation yields the following tokens:
$d_1$: [the, star, in, our, solar, system, provides, heat, and, light] \\
$d_2$: [that, hollywood, star, walked, the, red, carpet, for, the, movie, premiere] \\
$d_3$: [astronomers, observe, distant, stars, and, galaxies, using, telescopes]

\textbf{Step 2: Remove Stop Words} \\
Removing the specified stop words—\{the, in, our, and, that, for\}—leaves:
$d_1$: [star, solar, system, provides, heat, light] \\
$d_2$: [hollywood, star, walked, red, carpet, movie, premiere] \\
$d_3$: [astronomers, observe, distant, stars, galaxies, using, telescopes]

\textit{Note: "stars" in $d_3$ is treated as a distinct token from "star" since stemming was not requested.}

\textbf{Step 3: Construct the Inverted Index} \\
Mapping each unique term to the document IDs containing it:
\begin{itemize}
    \item astronomers: [$d_3$]
    \item carpet: [$d_2$]
    \item distant: [$d_3$]
    \item galaxies: [$d_3$]
    \item heat: [$d_1$]
    \item hollywood: [$d_2$]
    \item light: [$d_1$]
    \item movie: [$d_2$]
    \item observe: [$d_3$]
    \item premiere: [$d_2$]
    \item provides: [$d_1$]
    \item red: [$d_2$]
    \item solar: [$d_1$]
    \item star: [$d_1, d_2$]
    \item stars: [$d_3$]
    \item system: [$d_1$]
    \item telescopes: [$d_3$]
    \item using: [$d_3$]
    \item walked: [$d_2$]
\end{itemize}

\subsection*{2. TF-IDF Term-Document Matrix}
The total number of documents is $N = 3$. There are 19 unique terms.
Since each term appears exactly once in any document it belongs to, Term Frequency ($TF$) is simply 1 if the term is present and 0 otherwise.

For Inverse Document Frequency ($IDF$), we use:
\[ IDF = \log_{10}(N / df) \]
\begin{itemize}
    \item The token "star" appears in 2 documents ($d_1$ and $d_2$), so $df=2$:
    \[ IDF(\text{star}) = \log_{10}(3 / 2) = \log_{10}(1.5) \approx 0.176 \]
    \item All other 18 terms appear in exactly 1 document ($df=1$):
    \[ IDF(\text{other\_terms}) = \log_{10}(3 / 1) = \log_{10}(3) \approx 0.477 \]
\end{itemize}

The final term weights are calculated as:
\[ w = TF \times IDF \]

\begin{table}[htbp]
\centering
\begin{tabular}{|l|c|c|c|c|l|}
\hline
\textbf{Term} & \textbf{df} & \textbf{IDF} & \textbf{$w$ in $d_1$} & \textbf{$w$ in $d_2$} & \textbf{$w$ in $d_3$} \\
\hline
astronomers & 1 & 0.477 & 0 & 0 & 0.477 \\
carpet      & 1 & 0.477 & 0 & 0.477 & 0 \\
distant     & 1 & 0.477 & 0 & 0 & 0.477 \\
galaxies    & 1 & 0.477 & 0 & 0 & 0.477 \\
heat        & 1 & 0.477 & 0.477 & 0 & 0 \\
hollywood   & 1 & 0.477 & 0 & 0.477 & 0 \\
light       & 1 & 0.477 & 0.477 & 0 & 0 \\
movie       & 1 & 0.477 & 0 & 0.477 & 0 \\
observe     & 1 & 0.477 & 0 & 0 & 0.477 \\
premiere    & 1 & 0.477 & 0 & 0.477 & 0 \\
provides    & 1 & 0.477 & 0.477 & 0 & 0 \\
red         & 1 & 0.477 & 0 & 0.477 & 0 \\
solar       & 1 & 0.477 & 0.477 & 0 & 0 \\
star        & 2 & 0.176 & 0.176 & 0.176 & 0 \\
stars       & 1 & 0.477 & 0 & 0 & 0.477 \\
system      & 1 & 0.477 & 0.477 & 0 & 0 \\
telescopes  & 1 & 0.477 & 0 & 0 & 0.477 \\
using       & 1 & 0.477 & 0 & 0 & 0.477 \\
walked      & 1 & 0.477 & 0 & 0.477 & 0 \\
\hline
\end{tabular}
\end{table}

\subsection*{3. Boolean Retrieval}
\textbf{Query:} ``star light''

\textbf{Retrieving documents using the inverted index:}
From the inverted index constructed in Question 1:
star $\rightarrow$ [$d_1, d_2$] \\
light $\rightarrow$ [$d_1$]

\textbf{Which documents match the query terms?}
For a multi-term query like "star light", we usually use the Boolean AND logic. This means the document has to be present in the posting lists of both words. So we take the intersection of the list for "star" [$d_1, d_2$] and the list for "light" [$d_1$]. The common document is just $d_1$. Therefore, $d_1$ is our matching document. If we used OR logic instead, we would take the union and get both $d_1$ and $d_2$, but AND is much more common for this type of search.

\subsection*{4. Cosine Similarity and Ranking}
\textbf{Representing the query using TF-IDF:} 
The query vector $q$ contains two terms: "star" and "light" ($TF=1$ for both).
Using the previously computed IDF values, the weights are:
\begin{itemize}
    \item $w_{\text{star}, q} = 1 \times 0.176 = 0.176$
    \item $w_{\text{light}, q} = 1 \times 0.477 = 0.477$
\end{itemize}

Query vector length is:
\[ ||q|| = \sqrt{0.176^2 + 0.477^2} = \sqrt{0.030976 + 0.227529} = \sqrt{0.258505} \approx 0.508 \]

\textbf{Compute cosine similarity between query and each document:}
\begin{itemize}
    \item \textbf{Document $d_1$:} Contains 6 terms ("star" weight $0.176$, five terms weight $0.477$).
    \[ ||d_1|| = \sqrt{0.176^2 + 5 \times (0.477)^2} \approx \sqrt{0.030976 + 1.137645} \approx 1.081 \]
    \[ q \cdot d_1 = (0.176 \times 0.176) + (0.477 \times 0.477) = 0.030976 + 0.227529 = 0.258505 \]
    \[ \text{Cosine}(q, d_1) = \frac{0.258505}{0.508 \times 1.081} = \frac{0.258505}{0.549148} \approx 0.470 \]
    
    \item \textbf{Document $d_2$:} Contains 7 terms ("star" weight $0.176$, six terms weight $0.477$).
    \[ ||d_2|| = \sqrt{0.176^2 + 6 \times (0.477)^2} \approx \sqrt{0.030976 + 1.365174} \approx 1.181 \]
    \[ q \cdot d_2 = (0.176 \times 0.176) + (0 \times 0.477) = 0.030976 \]
    \[ \text{Cosine}(q, d_2) = \frac{0.030976}{0.508 \times 1.181} = \frac{0.030976}{0.599948} \approx 0.051 \]

    \item \textbf{Document $d_3$:} $d_3$ does not contain "star" or "light".
    \[ q \cdot d_3 = 0 \]
    \[ \text{Cosine}(q, d_3) = 0 \]
\end{itemize}

\textbf{Final Ranking:}
\begin{itemize}
    \item 1st: $d_1$ (score $\approx$ $0.470$)
    \item 2nd: $d_2$ (score $\approx$ $0.051$)
    \item 3rd: $d_3$ (score = $0$)
\end{itemize}

\textbf{Is the ranking desirable? Justify:}
Yes, this ranking makes a lot of sense. The query is "star light", which is related to space and astronomy. Document $d_1$ is talking about the sun giving heat and light, so it is the most relevant one and gets the highest score. On the other hand, document $d_2$ is about a movie actor (a "Hollywood star"), which has nothing to do with actual light. So $d_2$ getting a very low score is exactly what we want.

\subsection*{5. Analysis of Word Sense Ambiguity}
Consider the query: ``movie star''

\textbf{Which document(s) should ideally be retrieved?}
Ideally, for the query "movie star", we only want document $d_2$ because it is actually talking about an actor at a movie premiere.

\textbf{Does the system retrieve the correct document(s)?}
However, the system will retrieve both $d_1$ and $d_2$. Since both documents contain the word "star", they will both get some non-zero similarity score with the query.

\textbf{Explain how word sense ambiguity ("star") affects retrieval:}
This happens because the word "star" has multiple meanings (polysemy). In $d_1$, it means a real star in the sky, but in $d_2$ it means a famous actor. The problem with the basic TF-IDF Vector Space Model is that it only looks at exact word matches. It doesn't understand the context or meaning behind the words. Because it treats the space star and the actor star as the exact same word, it gives a higher score to $d_1$ even though it is completely irrelevant to movies. This ruins the precision of our search results.

\section*{Part 2: Building and Extending an IR System}
\textbf{1. Implement the retrieval component using TF-IDF} \\
The retrieval component has been fully implemented in the provided Python template utilizing the Term Frequency-Inverse Document Frequency (TF-IDF) vector space model. 

\begin{figure}[htbp]
    \centering
    \includegraphics[width=0.95\textwidth]{code1.png}
    \caption{Retrieval Component Implementation (Part 1)}
\end{figure}

\begin{figure}[htbp]
    \centering
    \includegraphics[width=0.95\textwidth]{code2.png}
    \caption{Retrieval Component Implementation (Part 2)}
\end{figure}

\begin{figure}[htbp]
    \centering
    \includegraphics[width=0.95\textwidth]{code3.png}
    \caption{Retrieval Component Implementation (Part 3)}
\end{figure}


\section*{Part 3: Evaluating your IR System}
\textbf{1. \& 2. Evaluation Measures and Global Metrics} \\
Using the evaluation template, the system was run on the Cranfield dataset. The averaged metrics across all queries at $k=10$ are:
\begin{itemize}
    \item \textbf{Precision@10:} 0.2507
    \item \textbf{Recall@10:} 0.3710
    \item \textbf{F$_{0.5}$-score@10:} 0.2568
    \item \textbf{Average Precision / MAP@10:} 0.2691
    \item \textbf{nDCG@10:} 0.4203
\end{itemize}
\textit{Note: The standard script didn't print out Mean Reciprocal Rank (MRR), but MAP came out to 0.2691 at $k=10$.}

\textbf{3. Evaluation Metric Plots and Observations} \\
\textit{[Placeholder: Insert graphs of all evaluation metrics vs k (e.g., eval\_plot\_part3.png) here]}

\textbf{Observations based on the plots:}
\begin{itemize}
    \item \textbf{Precision vs. k:} We can see that precision is highest at $k=1$ (0.6000) and slowly goes down as we increase $k$, reaching 0.2507 at $k=10$. This means our system is putting the most relevant documents at the very top. But as we look at more documents (larger $k$), we start bringing in more irrelevant ones, so precision drops.
    \item \textbf{Recall vs. k:} Recall does the exact opposite of precision. It starts low at 0.1048 for $k=1$ and keeps going up, reaching 0.3710 at $k=10$. This makes sense because as we retrieve more documents, we are bound to find more of the total relevant documents that exist in the dataset.
    \item \textbf{Stabilization:} The nDCG score drops a bit at first but then stays pretty steady around 0.42 after $k=5$. This shows that the quality of our ranking for the top few results is fairly stable compared to the ideal perfect ranking.
\end{itemize}

\textbf{4. Run Time Execution} \\
I used the standard Python \texttt{time} module to measure the execution time. For the entire IR process—which includes loading the dataset, preprocessing the text, building the index, running all 225 queries, and evaluating them—the total run time was roughly \textbf{2.19 seconds} for the 1400 documents.


\section*{Part 4: Analysis}
\textbf{1. Limitations of the Vector Space Model (VSM)} \\
The main issue with the Vector Space Model is that it only does exact word matching (the "Bag of Words" approach). It completely ignores grammar, word meanings, and synonyms, which causes a lot of problems.

\textbf{Concrete Example from Cranfield Dataset:}
\begin{itemize}
    \item \textbf{Query 1:} "what similarity laws must be obeyed when constructing aeroelastic models of heated high speed aircraft ."
    \item \textbf{Relevant Document (Doc 29):} "... determination of transient temperatures and thermal stresses in simple models intended to simulate parts or the whole of an aircraft structure of the built-up variety subjected to aerodynamic heating ..."
\end{itemize}
If we look at the provided ground truth (\textit{cran\_qrels.json}), Document 29 is actually very relevant to this query. But our TF-IDF system ranked it all the way down at \textbf{109}. The reason it failed so badly is because the query uses the specific word "heated", but the document uses similar words like "thermal", "temperatures", and "heating" instead. Since the VSM treats every unique word as a completely different thing, it doesn't know that "heated" and "heating" mean the same thing. So, the cosine similarity ends up being very low.

\textbf{2. The Zero-Result Problem (Out-of-Vocabulary Issue)} \\
\textbf{What happens in theory:} If we give the system a query where none of the words exist in our vocabulary (Out-Of-Vocabulary or OOV terms), the Term Frequency ($TF$) for those words will be zero. Because of this, the entire query vector just becomes an array of zeros.

\textbf{Effect on the retrieval results:} To find the cosine similarity, we have to calculate the dot product between the query vector and the document vectors. But since the query vector is all zeros, every single dot product will just evaluate to $0.0$. This means every document gets the exact same score of zero. The system has no way to rank them, so it just returns the documents in the exact order they were loaded into the program.

\textbf{Example Query showing this problem:}
\begin{itemize}
    \item \textbf{Query 2:} "intergalactic starships flippity floppity"
\end{itemize}
I tested the system with these completely random OOV words. Since none of these words are in the Cranfield documents, the code just outputted the very first five documents it had in memory: \texttt{[1, 2, 3, 4, 5]}. It didn't actually do any real matching because it couldn't find any of the words.

\end{document}
